\documentclass[UTF8,a4paper,11pt]{ctexart}
\usepackage[margin=24mm]{geometry}
\usepackage{amsmath,amssymb,booktabs,longtable,array}
\usepackage{xurl}
\usepackage[colorlinks=true,linkcolor=blue,urlcolor=blue,citecolor=blue]{hyperref}
\hypersetup{pdftitle={赛题19：动态算子图编译与并行调度研究报告},pdfauthor={Hackathon项目组}}
\setlength{\parskip}{4pt}
\setlength{\emergencystretch}{3em}
\Urlmuskip=0mu plus 1mu
\title{赛题19：动态算子图编译与并行调度\\\large 参数化块任务图、源码浮点语义与独立验证}
\author{本地项目：\texttt{19-dynamic-graph}}
\date{2026年10月11日}
\begin{document}
\maketitle
\begin{abstract}
本项目将受支持的实际C源码降低为包含可执行字节码和符号循环、访存表达式的Metadata，面对动态矩阵规模执行Cholesky分解。不能证明可变换的源码由通用VM保持源码顺序；能够证明完整运算图的左看、官网右看及连续分块形式，使用参数化POTRF/TRSM/UPDATE任务DAG。优化保持每个输出位置的乘加归约次序，禁止隐式FMA；显式FMA作为独立源码语义保留。

本轮已完成任务DAG、因子块打包、闭式依赖后继、独立累加器与CPU检测SIMD，并扩大源码前端、helper顺序求值、计数型while证明和整数转换测试。一般稠密SPD的最大规模$N=10000$已经运行原始C参考程序，并对完整800,000,008字节输出逐字节比较通过。报告区分一般稠密工作负载与对角/带状结构专门化，区分计算、任务图、文件IO与外部进程时间，并测量4/8/16线程伸缩和实际RAM峰值。保守row默认与显式memory配置各7项CTest通过；owned输入复用令一般dense10000峰值约2.02GB降至1.22GB，保留吞吐负结果。所有成绩是公开本地测试证据，不代表官方隐藏评测分数。
\end{abstract}

\section{项目目标与题面契约}
官网题目要求从朴素或分块Cholesky的C源生成算子符号依赖Metadata，再由Runtime处理动态Shape，$3\leq N\leq10000$。题面分数由Metadata大小15分和计算耗时85分构成，IO不计入计算性能，结果完全一致为硬门槛\cite{official}。本项目因此同时报告文件字节数、计算时间与严格结果比较；没有正式评分函数、评测CPU或官方数据包时，不推算最终分数。

官网源码图片采用右看更新：先对角开方，再计算本列下方因子，随后将乘积累加到后续下三角工作区。图片没有给出\texttt{div}/\texttt{madd}定义或输出初始化。明确的本地契约为：$\operatorname{div}(x,y)=x/y$，非融合$\operatorname{madd}(x,y,z)=x\times y+z$，$L$初始全为正零，上三角保留正零；另设显式\texttt{fma}源码验证不同语义。自设分块C不是官方分块源码。

输入、输出均采用小端\texttt{uint32}矩阵数$M$，随后每矩阵一个\texttt{uint32}维度$N$与行优先$N^2$个FP64值。单矩阵文件长度为$8+8N^2$字节，多矩阵长度为$4+\sum_z(4+8N_z^2)$。本地额外限制$0\leq M\leq100000$，允许空批次；这是资源保护契约，官网未给出该上限。输入为有限SPD矩阵，默认IEEE浮点环境，舍入到最近值；不改变FTZ/DAZ或舍入模式。符号零属于严格比较内容。

\section{重要工作及其对本项目的影响}
\subsection{LLVM与多面体分析}
Polly在LLVM IR上识别可建模区域，构建计算与访存模型，变换后重新生成IR；其文档还说明规范化、内联及LLVM优化管线位置会影响可分析性\cite{polly}。这支持本项目先建立符号循环与访存表达式、再进行可证明变换的分层设计。目前实际前端是Clang AST降低，并导出LLVM IR供审计，没有集成Polly或实现任意C的多面体编译器。整数地址交换、等价循环边界的规范化只用于整数符号证明，浮点归约表达式不作代数重排。

\subsection{任务化线性代数}
Buttari等的分块线性代数工作将Cholesky拆成对角分解、三角求解和尾部更新任务，通过依赖DAG释放并行度，并强调数据分块与任务调度共同影响局部性\cite{buttari}。LAPACK的\texttt{DPOTRF}采用阻塞形式与Level-3 BLAS\cite{lapack}。本项目采用相似的任务分解，但保留官网的独立$A$与$L$累加工作区，以及逐位置正向累加顺序；直接替换为普通BLAS不能自动保证给定源码的逐位真值。

PaRSEC以符号任务流表示依赖，其JDF可以在运行时局部计算依赖而不展开整个DAG\cite{parsec}。本项目先把堆分配后继列表改为闭式后继公式，再实现每目标块一个当前panel状态，将研究候选的常驻状态从$O(t^3)$降至$O(t^2)$，其中$t=\lceil N/B\rceil$。这只减少状态，不减少必须执行的浮点任务；一般符号图、工作窃取和局部性调度仍有工作余量。StarPU的codelet、数据句柄和读写访问模式能帮助运行时自动建立依赖\cite{starpu}；本项目借鉴其依赖与kernel分离的思路，未引入StarPU，也未测试GPU执行。

\subsection{确定性FP64与源码一致性}
LLVM的浮点语义区分默认环境、严格环境及fast-math标记；\texttt{reassoc}和\texttt{contract}允许的变换会影响归约与融合\cite{llvmfp}。Clang文档说明，仅指定\texttt{-fno-fast-math}并不能取代显式\texttt{-ffp-contract=off}\cite{clangfp}。因此参考C和Runtime均同时关闭fast-math与隐式收缩。

Ahrens、Demmel和Nguyen通过可复现累加器使求和结果与求和顺序无关\cite{repro}。本项目的要求更具体：与给定普通逐项C运算的位模式相同。由此推得，改用更准确或顺序无关的累加器也不能自动满足源码真值；本项目保留每个位置的加法链，而只并行不同位置。显式FMA源使用FMA语义，非融合源始终先乘后加。

\section{算法与实现}
\subsection{从实际源到可执行Metadata}
编译链为实际C源、Clang JSON AST、类型与声明身份、栈字节码及符号访存IR、受信Metadata、动态Shape Runtime。变量名、函数名和源hash不参与算法分派。改变实际乘加算式会改变程序，并使不满足证明条件的源码回退到VM；参考结果由另一条路径直接编译原始C源产生。

支持签名为\texttt{int/void f(double *A,double *L,int n[,int block])}，输入参数亦可为\texttt{const double*}。支持有符号32位\texttt{int}、\texttt{double}、直接数组下标、\texttt{for/while/if}、条件与短路表达式、基本算术、\texttt{sqrt/fabs/fmin/fmax/fma}，以及纯标量helper。新增helper可以包含按声明顺序初始化的不可变SSA局部值，最终一个return；有副作用、控制流或未初始化读取的helper拒绝。

VM中的整数算术结果和double转int均进行int32范围检查，零规范为正零。该处理修复了\texttt{trunc(-0.01)}产生负零、而C转换到int后再提升应为正零的差异；整数除法、余数、一元负号也有相应回归。未定义的整数溢出或除法不作为可执行C语义接受。

Metadata保留完整可执行程序，以及循环、父循环、读写数组和参数化下标。重复符号表达式编码为共享表达式DAG，可无损还原；可选expanded模式用于审计。来源与Metadata的SHA记录在分析JSON中，可按需内嵌来源信息。程序文件与矩阵维度无关；运行时绑定$N$及源码的可选\texttt{block}参数。物理调度块$B$是另一个参数，不要求等于源码分块参数。

\subsection{参数化块任务图}
令$t=\lceil N/B\rceil$，对$0\leq k<t$定义对角任务$P(k)$，对$k<m<t$定义三角求解$T(m,k)$，对$k<h\leq m<t$定义尾部更新$U(m,h,k)$。依赖为
\begin{align*}
P(k)&\leftarrow U(k,k,k-1),\\
T(m,k)&\leftarrow P(k),\ U(m,k,k-1),\\
U(m,h,k)&\leftarrow T(m,k),\ T(h,k),\ U(m,h,k-1).
\end{align*}
域外依赖省略；对角更新中的同一$T$只计一次。任务总数为
\[
Q(t)=\sum_{q=1}^{t}\frac{q(q+1)}2=\frac{t(t+1)(t+2)}6.
\]
每个任务使用8字节：三个16位块坐标、任务种类和未完成依赖数。节点后继由闭式索引计算；不为每个任务分配\texttt{vector}边列表。就绪优先队列偏向靠近对角的目标列，工作线程持续存在到整个矩阵完成。更新完成后才释放后继；每个输出块对不同$k$的更新形成严格链。

任务$P$处理对角块的原始递增列次序，$T$处理对应行块的原始除法和块内累加，$U$按递增源列更新尾部块。同一输出位置经历的乘法、加法、减法、除法和开方次序不变；不同位置的独立运算可异步。循环尾部统一以真实$N$截断，支持奇数维度、$B>N$与非整块。

\subsection{缓存、独立累加器与CPU SIMD}
已经完成的因子块按块内列打包，提供连续的右侧因子访问。新增缓存为
\[
8B^2\frac{t(t+1)}2\ \text{字节}\approx4N^2,
\]
输入和输出合计为$16N^2$字节。微内核同时维护4、8或16个\emph{不同输出位置}的累加器，不拆分单个位置的归约树。因子数据不可变，更新块之间的写区域独立。

AVX2路径放在单独编译单元，通用CPU检测器使用基础指令集，检测AVX2与FMA特性后才调用。该方法遵循GCC关于CPU检测和架构专用代码分离的说明\cite{gccx86}。非融合路径使用分别的向量乘法和加法；只有显式FMA源码使用\texttt{fmadd}。\texttt{--simd portable}强制回退。本报告不把跨CPU、不同数学库或不同浮点环境的逐位一致性作为已验证结论。

\subsection{IO和失败处理}
Runtime逐矩阵读取、计算、写出，内存中不保留整批$M$个矩阵。检查维度、长度、尾部数据、Metadata参数和跳转；错误结果不发布，临时\texttt{.partial}文件清理。为避免手工设置极小$B$展开不可负担的DAG，最多物化1600万个任务；$N=10000,B=32$处于此限额内。

\section{验证方法与已完成成果}
\subsection{独立参考与数据集}
参考工具直接包含并编译实际C源，使用GCC15.2、\path{-O2}、\path{-fno-fast-math}与\path{-ffp-contract=off}。它不使用Runtime、AST VM或复制Runtime的算法产生真值。比较覆盖头部、全部下三角与上三角、所有符号零；不以残差或容差替代逐字节比较。

一般稠密SPD用对称、严格对角占优矩阵构造：对角$N+2$，非对角为确定性哈希产生的正负小数。FP压力数据包括$DCD$混合尺度、真实subnormal元素、约$2^{900}$量级大数、接近秩一SPD和正负零。自设源码覆盖左看、官网右看、动态连续分块、显式FMA、变量与循环形式变化、SSA helper、改变算式以及通用while/if程序。

\begingroup\small\setlength{\tabcolsep}{3pt}
\begin{longtable}{@{}p{40mm}rrp{56mm}@{}}
\caption{本轮里程碑；各批证据独立保存，不合并成官方测试数量}\\\toprule
批次&数值对照&性能执行&证据范围\\\midrule\endfirsthead
\toprule 批次&数值对照&性能执行&证据范围\\\midrule\endhead
块DAG与打包&60&240&4种源、奇数/尾块、$N=256\ldots2048$，5轮\\
闭式图与大规模&96&62&真实subnormal等，$N=4096,10000$原始C真值\\
8/16独立累加器&128&180&4种源FP压力，3个规模、5轮轮换\\
扩展类型前端&188&已计入对照&13种源；整数转换、负号、除法/余数\\
CPU检测SIMD&128&112&portable/AVX2、4种源、3个稠密规模\\
对角因子缓存&128&176&原始C严格比对、5轮A/B与最大规模\\
Metadata紧凑编码&24&不作性能评分&4编码、3种源、VM和块DAG\\
对齐与就绪队列&136&264&奇数尾块、64/128列mask边界、10k重复\\
两行复用与线程唤醒&128&252&逐位置k链保持；5轮或最大3轮\\
自动物理块策略&16&222&4源、127/128与1024/1025切换、手动覆盖\\
四行向量复用&128&168&4源FP压力、大规模三轮及小规模负结果\\
线程相关块成本&0&200&同一候选、10个稠密规模、1/4线程\\
helper顺序与证明&49&不作性能评分&嵌套/prototype/线性SSA；另17项拒绝\\
带状与新稠密族&380&46&符号零、下溢、带宽边界、真实10k C\\
八行实验&128&128&强制dense；保留性能负结果\\
线程伸缩与RAM&0&48&4/8/16线程；普通10k及M3流式\\
线程与块几何&0&81&同一候选；B64/96/128/192，三轮\\
计数型while证明&63&不作性能评分&6种源；分块1/7/32/100；改算式回退\\
线程Shape分派&38&不作性能评分&4源、1536/1537及混合Shape边界\\
P/T独立位置SIMD&128&200&4源FP压力；B32负结果与B64/96收益\\
完整线程策略&144&116&自动B、1/4/8/16、M3与带状分别记录\\
有界解锁通知&208&156&相同运算与B；16线程负结果保留\\
整数比较选择helper&105&不作性能评分&分块for/while、FMA；另12项拒绝\\
缓存初始化与解锁通知&208&66&NaN poison、不清零；首次触页成本计入\\
连续列块工作区&209&94&正零、尾块、1/4/16；导出计入compute\\
混合块布局与并行导出&418&188&四构建轮换；保持负值与共享长尾\\
惰性当前panel任务图&418&104&heap/bucket；另30项异常退出\\
等待计数与行前缀&418&156&自动与B32；另30项异常退出\\
单worker闭式遍历&642&76&基础418+单worker224；另30失败\\
同源向量扫描&545&60&另288非有限拒绝、308880位/保护页单测\\
输入owned复用&518&80&另36失败、API所有权/const验证\\
auto复用与最终配置&854&74&另288拒绝+42失败；脚本失败/恢复保留\\\bottomrule
\end{longtable}\endgroup

\subsection{第一批缓存与调度收益}
表\ref{tab:phase1}是同批轮换A/B的5次中位数，旧基线为持续线程池、每列两阶段屏障，新版本为块DAG、打包与4独立累加器，均为普通稠密SPD。单线程的明显提升表明缓存局部性和独立累加器贡献很大；不能把全部收益归因于多线程调度。
\begin{table}[htbp]\centering\small
\caption{块DAG第一批，$B=32$；单位ms，计算包含分配与调度}\label{tab:phase1}
\begin{tabular}{rrrrrr}\toprule
$N$&旧1线程&新1线程&旧4线程&新4线程&同4线程加速比\\\midrule
256&5.334&1.274&5.962&0.916&6.51\\
512&98.964&12.283&32.755&3.975&8.24\\
1024&780.862&75.320&175.843&20.558&8.55\\
2048&6461.530&515.697&1286.530&135.039&9.53\\\bottomrule
\end{tabular}
\end{table}
$N=2048$时新引擎单线程相对旧单线程约12.53倍，自己的4线程相对单线程约3.82倍；这两个比例代表不同因素。本批$B=32$优于64、96、128，后续ISA优化仍重新搜索块大小。

\subsection{闭式图的收益和限制}
$N=4096,B=32$有357,760个任务。闭式后继把图构造约55ms降至约6--8ms，静态任务元数据为2,863,104字节；总体计算改善为几个百分点，并非数量级。$N=10000,B=32$有5,159,805个任务，静态任务元数据41,280,944字节，旧版本图构造约918ms，新版本约60ms。就绪队列及其他容器占用另计；这些数字不是题目评分中的磁盘Metadata大小。

一般稠密$N=10000$的原始C参考单次为504.670秒，输入生成34.767秒。第二批8种配置均与其完整800,000,008字节输出相同，SHA256为
\begin{center}\scriptsize\texttt{4015abfccfe5851a9c959939ff61ee44acf0290fcbf6192cdfd59a3a217eef13}\end{center}
紧凑图、4独立累加器、$B=32$的1线程计算49.174秒，4线程13.743秒；这是每配置一次的实测，不称为稳定中位数。$B=64$的4线程14.982秒而旧版本13.195秒，显示大规模单次观测与并发负载有波动，不能仅以图构造减少推断所有配置总时间必然下降。

\subsection{独立累加器与SIMD}
独立累加器批次中，$N=2048,B=32$的1线程4/8/16路中位分别569.888、415.119、379.598ms；4线程分别145.635、119.390、109.237ms。较小$N=512/1024$的4线程8路略优于16路，选择必须考虑规模和调度成本。

CPU SIMD批次已通过128组严格FP压力对照及112次稠密执行。$N=2048,B=32$三轮4线程中位为portable 95.707ms、AVX2 65.599ms，约降低31.5\%；$B=64/96/128$对应AVX2为85.364/115.572/147.139ms。一般稠密$N=10000,B=64$的单次1线程portable/AVX2为30.578/14.523秒，4线程为8.710/4.192秒，完整C参考字节相同。该次4线程读入约209ms、写出约609ms，均另外计时。

\subsection{对角因子缓存与紧凑Metadata}
第二次缓存优化让对角分解和行块求解直接使用已打包的对角因子，并在算出因子时立即写入缓存，取消独立打包扫描。$N=2048,B=64$的同批5轮4线程中位从86.592降至36.443ms；$N=4096$从452.281降至227.748ms。每个位置的浮点链未改变，128组FP压力与176次性能执行全部通过原始C参考。

最大规模单次结果并未显示相同收益：$N=10000,B=64$的4线程从3.074变为3.312秒，1线程13.167/13.175秒基本持平。本批不把对角缓存概括为所有规模加速，后续独立轮换批次比较缓存对齐和就绪队列，并保留运行区间及并发负载限制。

符号表达式共享DAG、来源信息转移到分析JSON、无参数字节码省去冗余零操作数，均保持可执行程序和展开符号图完全相同。左看、官网右看和自设分块的Metadata分别从同语义expanded编码2680/3280/4445字节降到2250/2655/3739字节，减少16.0\%/19.1\%/15.9\%。三种源的四种编码均通过VM与块DAG的独立C参考，共24次对照；这只是原始文件大小，不代入未知官方评分公式。

\subsection{就绪队列、两行复用与自动物理块}
分桶就绪队列按目标列分组，使用位掩码定位非空列，并优先推进对角任务；它与优先堆使用同一依赖图。136组FP/边界对照包含物理$B=1$跨64/128列mask边界，264次稠密执行全部通过。$N=4096,B=64$的4线程对齐heap/bucket中位233.590/234.562ms，最大规模三轮4153.250/4092.860ms且区间重叠。因此heap保持默认，bucket作为实验选项；未声称分桶稳定加速。64字节因子缓存对齐本批在$N=4096,B=64$从254.362改善到233.590ms。

两行向量内核让两个独立输出行共享右侧因子向量，保持各自递增k归约。它仍使用八个向量累加器，每个同时覆盖四个独立输出位置。选择通知只在退休释放一个或多个任务时唤醒worker，完成与异常时唤醒全部，避免没有新增可执行工作的通知。128组压力和252次稠密执行严格通过。表\ref{tab:rows}为本批固定$B=64$的4线程中位，不和前批最好值混算。
\begin{table}[htbp]\centering\small
\caption{两行复用与选择通知，同批轮换；计算ms}\label{tab:rows}
\begin{tabular}{rrrrr}\toprule
$N$&重复数&原单行&两行&两行加选择通知\\\midrule
2048&5&40.729&35.150&32.779\\
4096&5&239.099&221.794&199.052\\
10000&3&4059.180&3173.570&2984.180\\\bottomrule
\end{tabular}
\end{table}
最大规模最终候选区间2977.990--3663.950ms。同批1线程只执行一次，原单行/两行/两行加通知为14.555/11.394/12.047秒；通知候选单次串行负观测保留，不用三轮4线程收益推出所有配置必然改善。

自动物理块在$N\leq1024$选择32、较大$N$选择64，用户可显式覆盖；这是根据本地成本的任务几何策略，不影响数值结果或根据Shape查表结果。四种源共16次混合Shape对照检查127/128执行器切换与1024/1025物理块切换，并在portable/AVX2、1/4线程下保持原始C字节。222次稠密执行同时比较自动策略和显式$B=32/64$。最大规模自动4线程三轮中位3231.590ms，显式64为3376.470ms、显式32为5005.800ms；自动与64的差异来自观测波动，两者运行相同物理任务几何。有效块大小最小值/最大值与最后矩阵值由Runtime另行记录，混合Shape不以请求参数冒充实际B。

\subsection{helper求值、四行复用与线程相关成本策略}
helper的重复SSA局部变量曾通过表达式代入展开，可能指数增长；未使用的参数或局部值也可能漏掉本应执行的整数检查。新前端对这类调用按参数、声明、return的顺序降低为临时槽，所有参数先压栈，再写callee槽，因此同一helper嵌套不会覆盖外层待用参数。Clang声明链中的前向prototype归并到真实定义身份。20/40层重复SSA分别降低为146/266条指令，线性而非指数增长；资源上限为4096槽、100,000指令，并限制符号表达式展开。

原始C的49次整文件对照涵盖重复SSA、同函数嵌套、多参数嵌套、条件return、前向prototype和带未使用double值的Cholesky。17项拒绝检查覆盖未使用整数参数/局部值的溢出、真正递归、非有限常量、float/unsigned/wide helper返回类型及不符double契约的数学函数签名；其中类似Cholesky但携带未使用整数检查的源码，在auto、stage、tiled及VM请求下均拒绝。优化证明保守排除有未使用求值、int局部值或一般int参数运算的helper；不能只比较返回表达式后忽略其余求值。下一小节所述纯整数比较选择器是明确的无溢出例外。

新增计数型while证明把相邻的单变量初始化声明、while条件和末尾递增，仅在证明视图中规范化为for。实际字节码和循环审计仍保留while控制流；全量运算、循环域、递增步长和访存证明仍必须通过。63次原C比较覆盖左看、右看、显式FMA、分块while，源码block为1/7/32/100以及奇数/尾块；改变步长或算式的源码保持VM。该批实测20.073秒，不宣称一般while自动并行。

\subsection{分块边界的纯整数选择helper}
一般int helper保守留VM会让常见的\texttt{min(panel+block,n)}分块源失去任务DAG。新增的例外严格限制为两个int参数、无局部声明、一个return、直接比较两个参数并直接返回其中一个；实际传入表达式的类型也必须为int。helper自身的比较与选择不会发生int32算术溢出，调用参数的运算仍保留在完整循环/访存证明中。没有放开一般整数helper或丢弃其额外求值。

105次原C比较验证min边界的分块for、显式FMA与分块while，源码block为1/7/32/100；generic min/max保持VM。带double实际参数的整数转换、加入int局部拷贝等形式仍保守VM。12项拒绝在auto/stage/tiled/vm四种请求下覆盖未使用溢出局部值、显式及隐式参数转换溢出，且不为未定义C行为生成oracle。该批34.318秒，随后原有helper49/17回归37.548秒。min分块Metadata3810字节、while形式3655字节；它们较普通自设分块3739字节略增的原因是明确参数槽与调用求值，不能为减小Metadata而省略转换。

四行乘八列的AVX2内核用八个向量累加器共享右因子，每个位置仍按递增k处理。128组FP压力与168次性能对照通过。固定$B=64$，4线程的两行/四行中位在$N=2048$为32.157/30.408ms，在4096为197.180/183.399ms，在10000为2876.000/2527.850ms；最大规模四行区间2514.800--2665.420ms。该批最大规模1线程仅一次，10.152/8.816秒，不能当稳定中位数。较小规模有负收益，例如$N=512,B=64$的4线程1.936/2.162ms。因此自适应候选在$N\geq2048$用四行，否则两行。

独立中间规模成本批次共200次严格稠密对照，比较同一自适应二进制的$B=32/64$。1线程在$N=1152$为18.421/21.683ms，1536为42.849/52.100ms；4线程对应8.025/8.000ms、17.820/16.748ms。由此候选的单线程小块界限放到1536，多个线程保留1024；2线程等未实测线程数属于策略外推，用户可显式覆盖物理块。本地策略不承诺在其他CPU或负载下最优。

八行乘四列的另一种布局将右因子载入共享给更多行，但增加广播数量。独立128组强制dense FP对照与128次性能执行通过，仍为负结果：$B=64$、4线程的自适应候选/八行内核在$N=2048$中位30.737/31.820ms，4096为184.018/192.115ms，10000三轮2322.460/2510.260ms。它作为编译配置实验保留，不选作默认；不能只因理论载入复用增加就宣称更快。

\subsection{数据驱动带状专门化与新稠密工作负载}
Netlib的\texttt{DPBTRF}专门处理SPD带矩阵并提供带状存储接口\cite{lapackband}。本项目仍接受并输出完整$N^2$布局，先检查下三角的确切非零半带宽$w$，仅在$w\leq128$且$4w<N$时启用带状路径；显式\texttt{--structure dense}或较小band-limit可以关闭它。结构由实际数据扫描确定，不由文件名、Shape或预置标签确定。

对行$i$取$f=\max(0,i-w)$。带内$(i,j)$仅跳过归约中$k<f$的前缀，其余$k=f,\ldots,j-1$全部执行，含带内实际零。被省略前缀的一个乘数为确切零、另一个有限，初始累加器正零，在本地默认舍入环境下保持正零；之后的乘加链与原C相同。带外因子保持输入符号零；上三角仍为正零。不能省略中途零项，因为FMA下溢后的负零可能被后续加零改变。带状运算为$O(Nw^2)$，扫描、输出分配及完整格式仍为$O(N^2)$；当前带状kernel串行，指标明确记录实际worker为1。

结构化里程碑380次原始C整文件对照覆盖左看、官网右看、自设分块、显式FMA，normal/subnormal/large/tiny下溢与符号零、奇数规模，以及带宽127/128/129的分派边界；46次性能执行另外记录。稠密Toeplitz相关矩阵和正对角加三个外积的low-rank族，覆盖2047/2048/2049的向量策略边界，均保持各自原始C字节。它们未作为带状快速路径结果混入一般稠密基准。

真实$N=10000,w=7$带状输入为800,000,008字节，原始官网C的独立编译参考实际执行481.611秒，输入生成2.626秒。完整输出SHA256为\path{fb50d92b7bb6519d52bb69afba48b36745ac56f9bf7cf3e109adbcb53154a0c2}。同批三轮旧稠密内核/带状路径compute中位3002.110/184.861ms，外部wall中位3877.539/1211.165ms；带状compute区间172.044--215.752ms。其结构扫描30.324--33.211ms，kernel37.618--39.199ms，读入196.016--208.742ms、写出598.288--655.389ms。该收益只适用于已验证带结构，不能替代一般稠密最大规模结果。

一般稠密的早退结构扫描在同批$N=4096$的5轮中位226.439/233.907ms呈负观测，$N=2048$为33.721/32.822ms；没有证据说明所有规模都因此加速。新增路径保留可关闭选项，性能波动与其他项目并发限制继续适用。

\subsection{4线程以上伸缩与真实RAM峰值}
复用并核验已有普通稠密原C参考，48次伸缩执行均保持完整字节。固定$B=64$、同一二进制，表\ref{tab:scaling}为同批轮换中位；2048/4096各5轮，10000及混合流各3轮。4到16线程在普通最大规模约1.56倍，并非线性；混合流的16线程中位比8线程略慢，不能推广为任何批次更多线程必然更快。
\begin{table}[htbp]\centering\small
\caption{伸缩批次计算中位ms，固定B64}\label{tab:scaling}
\begin{tabular}{lrrr}\toprule
工作负载&4线程&8线程&16线程\\\midrule
普通dense2048&38.857&27.246&25.222\\
普通dense4096&194.826&132.536&120.677\\
普通dense10000&2535.090&1801.070&1630.260\\
M3:4096/10000/4096&3380.520&2227.910&2272.450\\\bottomrule
\end{tabular}
\end{table}
一般最大规模的4/8/16线程区间为2506.240--2593.880、1704.960--1900.300、1628.130--1657.110ms；wall中位3549.484/3056.392/2567.551ms。IO及分配占比随着计算缩短而增大。

通过Windows\texttt{GetProcessMemoryInfo}读取系统维护的\texttt{PeakWorkingSetSize}与\texttt{PeakPagefileUsage}，后者是进程峰值commit而非实际换页量\cite{winmemory}。PrivateUsage及宿主可用物理RAM另以20ms间隔采样。普通10000峰值工作集约1924.5MiB，峰值commit约1.879GiB；混合M3完整文件1,068,435,472字节，峰值工作集约1924.7MiB，和单个最大矩阵相近，支持逐矩阵流式内存不会按M累加的结论。宿主最小可用RAM约6.45GiB，但受其他项目共享负载影响，不归因于19号进程。资源采样会增加外部wall开销，不计入native compute。

随后81次同一二进制的线程/几何批次比较4线程B64与8/16线程B64/96/128/192。任务数量在10000从B64的657,359降到B96的198,485、B128的85,320、B192的26,235，减少退休次数确实有潜力，但更大的串行P/T和缓存工作集抵消收益。8线程B64/96/128/192中位1945.690/1932.630/2334.030/2400.190ms；16线程2372.620/1962.650/1878.700/2289.820ms，范围大且部分重叠，和前批16线程1630.260ms也有明显共享负载波动。该批不足以单独确定默认；后续将P/T内独立位置向量化和高线程数B96组合为完整策略验证。

\subsection{P/T独立位置SIMD的收益与负结果}
P/T块内的每次更新都对不同输出列执行同一个已确定因子的乘加。广播该因子并向量化这些独立位置，不改变任一位置的源列顺序，也不把串行归约改成树。128次强制dense浮点压力与200次性能执行均经原始C全字节校验。B64的2048单线程中位96.692/89.824ms；B96的4096单线程644.804/604.312ms，区间632.843--657.380与594.514--611.965ms；其4线程183.591/173.666ms，区间177.977--194.522与170.364--177.117ms。这些中等规模支持P/T内SIMD的收益。

B32的2048单线程93.945/94.866ms、4线程36.239/39.097ms为负观测。10000的单线程B64为8531.410/8730.380ms，仅一次；4线程B96为2835.370/2994.660ms亦负。16线程B96为1959.040/1910.230ms，但新区间1462--2638ms宽，不能称为稳定最大规模收益。由此下一完整策略仅对B至少64开启P/T SIMD，并将8/16线程的4096以上物理B候选设为96；仍需组合后的轮换与M3流式基准决定默认。

\subsection{完整线程策略验证}
随后完整候选把P/T SIMD限制到$B\geq64$，并在至少8个worker且$N\geq4096$时自动选择B96；1/4线程原32/64界限保留。144次浮点压力与116次严格性能执行通过，复用已核验C参考，整批303.960秒。表\ref{tab:thread-policy}比较旧Shape策略与完整候选的同批中位；10000单线程仅一次，其余每配置三次。
\begin{table}[htbp]\centering\small
\caption{完整策略计算ms；普通dense与M3流分别列出}\label{tab:thread-policy}
\begin{tabular}{@{}lrrr@{}}\toprule
工作负载&线程&旧策略&完整候选\\\midrule
dense2048&1&92.950&87.095\\
dense4096&1&715.265&627.143\\
dense4096&4&221.434&189.402\\
dense4096&8&136.119&128.316\\
dense4096&16&109.096&97.664\\
dense10000&1&9053.970&9417.710\\
dense10000&4&2678.300&2336.570\\
dense10000&8&1664.370&1700.280\\
dense10000&16&1648.560&1368.370\\
M3:4096/10000/4096&4&3282.370&3473.990\\
M3:4096/10000/4096&8&2578.390&2472.920\\
M3:4096/10000/4096&16&2505.220&2358.660\\\bottomrule
\end{tabular}\end{table}
4096的单线程12.3\%改善主要来自kernel而不是增加线程，其4线程14.5\%改善属于相同线程数下的内核/策略差异；16线程不能与单线程比率相乘成官方加速比。10000的16线程compute中位1648.560/1368.370ms，kernel1435.880/1201.510ms，graph41.169/40.404ms，wall2625.863/2272.141ms。任务数量减少没有自动带来同等比例graph时间减少，因graph还包含约400MB因子缓存初始化。

10000的8线程compute负观测1664.370/\allowbreak{}1700.280ms仍保留，即便其kernel1492.430/\allowbreak{}1331.020ms改善；扫描、分配与共享RAM压力可能影响compute的其他部分。4线程新区间2294.490--6424.710ms包含一轮明显长尾，未剔除；旧区间2432.280--2687.950ms。M3的4线程compute亦负，wall却5943.328/\allowbreak{}5169.361ms，说明IO/启动与compute口径不能混同。该批宿主最小可用RAM约3.576GiB，同时其他项目使用CPU/GPU；部分末尾观测还与下一候选构建重叠。进程自身峰值工作集2,018,201,600字节，约1.880GiB，并未按M累加。带状10000两版本均用同一串行结构kernel，其compute200.975/\allowbreak{}226.705ms为负观测，不应归因于未执行的P/T优化。此候选是可覆盖的本地成本启发式，未证明8/\allowbreak{}16线程在所有负载或机器上全局最优。

\subsection{有界唤醒的负结果}
对同一完整线程策略，实验把依赖状态更新后先解锁，再按\texttt{min(released,workers-1)}逐次notify\_one，完成与错误仍唤醒全部。谓词与依赖状态仍受同一mutex保护，208次FP/奇数/ISA对照及156性能运行均通过原C字节，整批489.754秒。但减少被唤醒线程数没有保证总成本下降：1024的4线程7.535/9.392ms，4096的16线程104.563/113.074ms，10000的16线程1405.130/1652.730ms均负；后者kernel1216.810/1459.080ms也负。

2048的4线程33.060/30.080ms、16线程23.362/20.842ms呈正观测，M3的4线程3784.760/3437.740ms及8线程2616.990/2330.750ms也正，但16线程2253.080/2351.250ms负。固定B32的2048在4/8/16线程35.025/44.028/64.330ms变成36.377/44.746/65.189ms，均负。该候选因此保留为默认关闭的实验。每次notify\_one本身有成本，批量释放多项时多次调用也可能比一次notify\_all慢；这个解释是结合实现的推测，并未测量操作系统唤醒事件。下一候选仅把原有通知方式移到锁外，不改变唤醒数量，另检验因子缓存清零成本。

\subsection{因子缓存初始化与原通知移到锁外}
该批进一步分离两个变量：unlock仅把原来的notify\_one/notify\_all选择移到锁外，lazy在unlock基础上取消因子缓存的冗余清零。这个缓存的所有实际读取都发生在P/T写入之后，P只读取本块已完成的前序位置，填充区不读取。输出L仍以正零初始化。另一个poison构建把所有缓存槽填为NaN；poison与lazy共208次原C比较覆盖4种源、subnormal/大数/近秩一/符号零、portable/ISA、奇数尾块、B大于N和小N的B1，全部通过，耗时38.130秒。此证据仅针对原行主序输出加因子缓存布局；后续把缓存改为累加工作区时，初始正零重新成为必要条件。

性能批次固定相同自动B策略，逐轮轮换三个构建，共66次严格核验，耗时173.605秒。表\ref{tab:lazy}保留全部分组中位；N=10000为3轮，混合M3为2轮，其余亦为3轮。
\begin{table}[htbp]\centering\small
\caption{缓存初始化实验；compute毫秒}\label{tab:lazy}
\begin{tabular}{rrrrr}\toprule
负载/线程&原策略&unlock&lazy&lazy相对原策略\\\midrule
2048/4&31.543&30.379&32.643&负\\
2048/16&18.694&19.922&17.206&正\\
4096/4&210.665&223.402&189.636&正\\
4096/16&106.929&110.283&100.042&正\\
10000/4&3196.950&2885.190&2692.640&正\\
10000/16&1337.160&1499.680&1485.640&负\\
M3/4&3748.155&3489.615&3456.605&正\\
M3/16&2353.970&2472.620&2056.330&正\\\bottomrule
\end{tabular}\end{table}
10000/4原策略区间2537.780--3271.420ms，lazy区间2423.210--3409.710ms，区间重叠，不声称稳定幅度；10000/16的负观测保留。lazy的graph中位在2048/4为0.040ms、4096/4为0.136ms、10000/4为2.378ms，原策略对应2.123/7.062/47.604ms，但首次触页与实际缓存写入迁入kernel，图初始化缩短不能单独视为端到端收益。10000/16的kernel原策略1133.260ms，lazy1324.450ms，反而增加；wall原策略2616.967ms，lazy2464.208ms，IO差异又改变了wall方向。实验时其他项目继续占用CPU/GPU，宿主最小可用RAM6,307,524,608字节，进程峰值工作集2,018,152,448字节，均作为共享环境观测。两个开关暂仍默认关闭，不以一次建图收益替换全流程选择。

\subsection{连续列块工作区的负结果与导出成本}
该候选把累加区与已完成因子合并为同一份下三角tile-major、块内column-major工作区。P/T重排独立位置而不重排每位置的k链，乘法左、右操作数顺序仍保留。累加工作区必须以正零初始化；结束后把已完成下三角导出到完整行主序L，上三角仍为正零。全L仍提前分配，因而峰值RAM没有降低；仅热工作集可能由原行主序L下三角加因子缓存约800MB变为约400MB。一个64乘64块为32KB连续存储，原L的大N行跨度可能分散在约64页；这只是地址几何推断，未采集TLB或缓存硬件计数器。

209项原C对照通过，耗时34.512秒；覆盖4种源码、FP压力、奇数和B大于N、B1、portable/ISA、1/4/16线程，源码block另有1/32/100。五项CTest也通过，174.410秒。随后94次性能对照168.838秒，在表\ref{tab:column}中明确计入串行导出。除N=10000单线程仅1轮、M3仅2轮外，各配置3轮。
\begin{table}[htbp]\centering\small\setlength{\tabcolsep}{3pt}
\caption{连续列块候选同批中位；全部毫秒}\label{tab:column}
\begin{tabular}{rrrrrr}\toprule
负载/线程&原compute&列块compute&原kernel&列块kernel&列块export\\\midrule
512/1&2.768&2.625&2.193&1.884&0.141\\
512/4&1.718&1.585&1.133&0.880&0.149\\
512/16&2.125&2.207&1.558&1.443&0.157\\
1024/1&15.801&14.440&13.522&11.703&0.707\\
1024/4&6.772&6.495&4.842&3.625&0.761\\
1024/16&10.094&11.004&7.966&8.004&0.850\\
2048/1&113.152&98.758&105.140&88.019&3.497\\
2048/4&31.458&32.734&23.181&21.251&3.095\\
2048/16&19.633&23.008&11.668&10.793&3.331\\
4096/1&819.546&702.240&792.013&661.752&14.836\\
4096/4&197.625&204.194&163.918&162.591&13.906\\
4096/16&113.715&114.066&86.162&73.410&13.373\\
10000/1&9115.610&12350.500&8869.010&12055.700&84.188\\
10000/4&3294.160&3549.350&2965.280&3153.380&98.103\\
10000/16&2265.810&2271.630&2055.420&1839.410&107.687\\
M3/4&3748.605&4445.390&3382.410&3998.485&145.844\\
M3/16&2120.265&2269.390&1821.425&1874.410&124.257\\\bottomrule
\end{tabular}\end{table}
不少kernel中位降低，但导出抵消了多线程收益；最大规模单线程kernel亦明显负。10000/4原compute区间2977.950--3621.750ms，列块3203.350--3694.340ms；10000/16原2101.500--2320.830ms，列块2172.080--2562.210ms。共享宿主同时仍运行其他CPU/GPU项目，最小可用RAM5,036,146,688字节，进程峰值工作集2,018,258,944字节，不能用不同批次的最小值推出因果。该候选默认关闭。检查实际AVX反汇编，热k循环没有向栈spill，故不把负结果归因为寄存器spill。下一候选将目标累加区改为块内行主序，复用原已验证的四行乘八列内核；P/T完成后用每worker小scratch改变为列因子布局，并独立检验最终位保持转置与并行导出。

\subsection{混合布局与位保持导出的独立对照}
混合候选令U任务的累加目标保持块内行主序，因子读取仍保持块内列主序，因而可复用已验证的原四行乘八列SIMD内核。P/T拥有目标块的独占写权限，用worker局部scratch顺序形成新因子；全部累加值消耗完毕之后、依赖退休之前，把目标块转换为列因子表示。消费者仅在完成后读，因此同一工作区随任务状态改变解释，不保留两份下三角大缓存。scratch惰性分配，每worker最多$8B^2$字节，B96且16worker理论上限1,179,648字节；仍提前分配完整L，不能宣称全流程峰值RAM降低。

导出另有默认关闭开关：4乘4 AVX2解包与lane置换仅搬运原FP64位，按独立块行分工；N不小于2048且worker多于1时并行，较小N保持串行。奇数尾块和对角下三角独立处理，上三角不写。混合串行与快速导出共418次原C压力对照通过，51.342秒。随后把原行主序、连续列块、混合串行、混合快速导出四个冻结构建轮换，188次性能执行完整通过，350.395秒。
\begin{table}[htbp]\centering\small\setlength{\tabcolsep}{3pt}
\caption{混合布局批次compute中位；毫秒}\label{tab:hybridtiles}
\begin{tabular}{rrrrr}\toprule
负载/线程&原行主序&连续列块&混合串行&混合快速导出\\\midrule
1024/1&17.128&14.521&13.807&14.375\\
1024/4&6.819&7.449&6.205&6.143\\
1024/16&9.935&10.519&10.753&10.516\\
2048/1&89.991&91.815&83.187&81.595\\
2048/4&29.563&30.967&30.125&28.596\\
2048/16&20.172&20.610&22.704&19.544\\
4096/1&700.708&845.260&591.991&662.861\\
4096/4&194.568&199.421&190.820&190.838\\
4096/16&106.007&107.886&118.323&108.243\\
10000/1&9227.100&11327.100&18891.800&9791.520\\
10000/4&3016.970&3483.620&3039.080&2886.730\\
10000/16&1630.900&1999.950&1875.280&2095.940\\
M3/4&3451.045&3546.595&5055.220&3143.960\\
M3/16&1634.995&1819.310&1878.425&1963.315\\\bottomrule
\end{tabular}\end{table}
导出阶段的明确局部收益为：4096/4串行14.018ms、快速6.350ms，4096/16为13.579/5.048ms，10000/4为95.365/33.105ms，10000/16为117.100/32.065ms。尽管导出减少，10000/16 compute仍负，不能以阶段改进代替端到端收益。10000/4原compute区间2817.680--3128.430ms，快速2870.760--3088.790ms重叠；16线程原1495.190--2270.770ms，快速1716.200--2199.220ms。M3/4串行混合有6660.510ms长尾，与另一轮3449.930ms一起保留；表中2轮中位是两数平均，不称为稳定收益。

该批早段512出现明显共享干扰：原单线程3.342/\allowbreak{}28.543/\allowbreak{}21.033ms，混合串行2.644/\allowbreak{}2.574/\allowbreak{}25.616ms，快速2.492/\allowbreak{}2.473/\allowbreak{}5.852ms。1024原单线程亦含129.252ms长尾。10000的混合串行和快速导出使用相同kernel代码，却单次kernel为18258.800/\allowbreak{}9510.960ms，故不能把差异归因于导出开关。全部原始观察留在JSON，未据此更新小N阈值。宿主最小可用RAM5,263,216,640字节，进程峰值工作集2,019,373,056字节。布局和导出开关暂仍默认关闭；下一批固定原行内核和shape，仅检验任务图状态压缩，并记录独立系统CPU计数帮助判断共享环境。

\subsection{当前panel惰性图：状态节省与性能边界}
候选不再为每个$(k,m,h)$物化一个状态，而为每个下三角目标块$(m,h)$保留一个8字节状态，依次执行$k=0,\ldots,h-1$的U，随后执行P或T。U结束后该状态推进一个panel；P/T结束后成为不可变因子。当前U就绪仍要求两个因子完成，当前T要求对角因子完成，先前U由状态推进隐含保证。工作线程持锁取出并复制任务，解锁仅执行计算，完成后重新持锁推进状态；就绪堆中的状态不修改。错误广播并等待所有worker退出，另检查无任务可执行且无active时的停滞，以及最终所有目标块已完成。

逻辑执行数仍为$Q=t(t+1)(t+2)/6$，常驻状态变为$S=t(t+1)/2$；1600万逻辑任务保护上限继续保留。它没有省略数学工作，也没有声称输出、因子缓存或总RAM降为$O(t^2)$。heap和bucket分别通过4种实际C源、非融合/显式FMA、subnormal、大数、近秩一、符号零、portable/ISA、奇数尾块与B大于N，总计418次完整字节比较，耗时67.036秒。另30项负pivot测试覆盖不同panel及1/4/16线程，均在限时内退出且没有完整或partial输出，耗时5.236秒。

\begin{table}[htbp]\centering\small
\caption{同shape的逻辑任务与常驻状态；字节仅为任务状态和panel索引}
\begin{tabular}{rrrrrr}\toprule
$N$&$B$&$Q$&$S$&旧字节&惰性字节\\\midrule
1024&32&5984&528&48128&4224\\
2048&64&5984&528&48128&4224\\
4096&64&45760&2080&366592&16640\\
10000&64&657359&12403&5260128&99224\\
10000&96&198485&5565&1588720&44520\\\bottomrule
\end{tabular}\end{table}

104次性能批次复用相同原C真值、源码、行主序kernel、通知与自动B策略，整批209.108秒。一般稠密4096的16线程compute中位旧/惰性为102.572/95.272ms，三轮范围99.912--109.238/94.732--99.642ms，kernel中位75.720/70.672ms；这个局部结果有价值，但其他配置未一致改善。10000的4/8/16线程compute中位分别为2844.350/2562.990、1570.820/1875.660、1628.270/1707.790ms，8和16线程为负。10000单线程仅一次9615.620/9178.440ms，不构成稳定结论；M3的4/8/16线程两轮中位为3584.335/3743.605、2702.365/2380.530、2133.900/2127.610ms。1024的1/4/8/16均为负，说明减少状态并不自动改善小任务开销。

该批4096第三轮4线程旧/惰性分别有1957.010/\allowbreak{}1752.040ms长尾，8线程977.520/\allowbreak{}988.571ms，同样保留。4线程两次外部wall为3000.990/\allowbreak{}2623.297ms，进程全线程CPU合计7156.250/\allowbreak{}6734.375ms，宿主其他CPU时间估计47218.750/\allowbreak{}30281.250ms。期间另一个项目进行32700 token GPU容量探测，主机与显存压力未隔离；这些计量只能描述共享负载，不能把尖峰归因为DAG实现。峰值工作集2,018,127,872字节，宿主最小可用RAM5,407,789,056字节。

\subsection{当前等待计数与因子行前缀：负结果亦保留}
第二候选缓存当前任务尚缺的1或2个因子通知，因子完成只减少相应当前panel的计数；目标推进后重新计算当前等待数，不缓存未来任意事件。每行T$(m,k)$依赖该行T$(m,k-1)$，最终P依赖最后T，因此完成因子的panel形成连续前缀。以每行一个int16前缀替代三角索引和随机Done状态读取，退休时检查严格递增不变量。比较$-1$初值时使用有符号类型，避免转换成size\_t。新增目录为$2t$字节，计入任务元数据；kernel、浮点表达式、自动B及通知策略保持不变。

heap/bucket 418次严格原C对照耗时59.524秒，另30项异常生命周期0.957秒。156次三候选性能耗时605.969秒，自动B与强制B32独立记录。表中S为原物化线程策略，L为惰性Done检查，P为缓存等待加行前缀；单位ms。普通最大规模单线程仅一次，其他单矩阵三次，M3两次。
\begingroup\small\setlength{\tabcolsep}{4pt}
\begin{longtable}{@{}lrrrr@{}}
\caption{同批compute中位：状态减少未形成全局稳定收益}\\\toprule
工作负载&线程&S&L&P\\\midrule\endfirsthead
\toprule 工作负载&线程&S&L&P\\\midrule\endhead
N2048&1&95.845&99.162&106.535\\
N2048&4&34.035&31.092&29.925\\
N2048&8&23.647&23.514&23.903\\
N2048&16&19.833&20.661&19.308\\
N4096&1&622.157&697.030&786.059\\
N4096&4&253.546&226.649&214.560\\
N4096&8&148.780&182.218&163.573\\
N4096&16&120.790&131.573&116.398\\
N10000&1&9415.250&9176.190&9410.580\\
N10000&4&3252.660&3110.070&3201.230\\
N10000&8&2214.760&2636.770&2789.540\\
N10000&16&2061.090&1785.840&1986.570\\
M3&4&3648.790&3545.940&3593.185\\
M3&8&2272.625&2788.170&2356.890\\
M3&16&2192.980&2604.815&3152.240\\
N4096/B32&4&362.082&314.891&314.839\\
N4096/B32&16&587.602&589.181&709.903\\
N10000/B32&4&4666.880&4595.630&4902.780\\
N10000/B32&16&10910.200&11812.000&12358.700\\
\bottomrule\end{longtable}\endgroup

P在2048的4线程局部改善，但在4096单线程、10000的8线程及M3的16线程明显为负。M3的16线程S/L/P进程全线程CPU中位为29796.875/35679.688/42445.313ms，宿主busy中位48.913\%/53.757\%/64.635\%；其他任务未隔离，不能据此证明具体因果。强制B32让10000逻辑任务增至5,159,805，16线程S/L/P中位10910.200/11812.000/12358.700ms，均远慢于自动B96。P该组三轮11139.200--26330.500ms，最大的26.331秒仍与完整原C结果逐字节相同，不删除为异常值。

该长尾外部wall 28428.794ms，input read 1235.830ms、graph 363.227ms、kernel 24997.700ms，进程CPU合计205375.000ms（kernel态123234.375ms），其他CPU时间估计386843.750ms。该时段RAM可用7,925,399,552字节，页错误491,894次；没有hard-fault计数，不能将其解释为磁盘分页，也不能用宿主负载独自解释全部调度差异。整批峰值工作集2,050,871,296字节，宿主最小可用RAM3,449,270,272字节。前缀模式的10000/B64任务元数据99,538字节，B96为44,730字节，B32为393,754字节；源Metadata仍为2655字节。O(t²)空间改善成立，时间结果混合，两个惰性图选项继续作为默认OFF的研究候选。

\subsection{单worker闭式拓扑遍历}
另一默认OFF候选\path{CHOL_SERIAL_TILE_WALK}只在实际worker为1时按目标列col、panel $k<col$、目标行row执行U，随后P(col)和本列T。此前列的因子均完成，每位置仍递增k；顺序对应heap的列/panel/行优先顺序。它不创建常驻Task或就绪队列，也不进入锁、条件变量；逻辑Q、初始化、实际数学任务和1600万逻辑cap保留。JSON记录\path{ready_queue=closed_serial}、常驻状态与ready peak均为0。多worker继续原物化DAG，不声称它们有算法改动。

基础heap/bucket418次原C对照63.918秒；额外224次全部single-worker auto/portable，覆盖四源、四FP模式、B1/7/31/64/96/128/1024及奇数/尾块，62.237秒，并断言真实closed遍历与零常驻状态。另30项错误生命周期5.498秒。76次性能306.705秒；每单矩阵三轮，含普通10000及强制B32的单线程三轮，M3两轮。表为旧S/候选W compute中位ms。
\begin{table}[htbp]\centering\small
\caption{单worker整版本对照；4线程为未改算法的控制观测}
\begin{tabular}{lrrrr}\toprule
工作负载&S单线程&W单线程&S四线程&W四线程\\\midrule
N512&2.735&2.599&1.793&1.862\\
N1024&15.805&15.421&6.485&6.797\\
N2048&96.382&90.506&32.563&31.145\\
N4096&632.523&809.327&197.345&196.032\\
N10000&9463.690&10204.800&2927.970&2778.400\\
\bottomrule\end{tabular}\end{table}

普通4096的S/W单线程中位632.523/809.327ms、10000为9463.690/10204.800ms，均为负。10000单线程范围S 9450.320--10210.100、W 10048.200--11316.600ms。强制B32的4096单线程S/W为983.809/927.363ms，10000为14554.300/14533.800ms基本持平；M3单线程两轮中位11330.600/11526.600ms。512、1024、2048有局部改善，未形成普通大规模稳定收益，因此默认仍OFF。四线程没有改变调度算法，却出现正负计时差异，说明编译实现与共享负载仍影响整版本对照，不能把所有差异归因为闭式遍历。

整批进程峰值工作集2,050,170,880字节，宿主最小可用RAM6,481,620,992字节。单worker常驻Task字节确实降为0，完整输入、输出和因子存储仍保持，未宣称总RAM消除。

\subsection{输入所有权复用：约800MB实测峰值节省}
在列块统一工作区、或行累加块转列因子的混合模式，P/T/U只读A并读写packed，不访问完整L。已有实现提前分配L，虽缓存仅保留下三角，峰值仍为A+L+packed。默认OFF的\path{CHOL_REUSE_INPUT_BUFFER}提供\path{execute_owned(Matrix input,...)}入口；CLI将读入的vector移动到该入口，不复制payload。只读\path{execute(const Matrix\&,...)}保持输入不变，不能复用输入。所有者输入的释放归execute\_owned，在成功返回或异常展开时完成；旧CLI变量原先直到写出结果后才释放，故新旧销毁计时不可直接混同。

第一阶段只对显式\path{--structure dense}、已证明native图且列块布局生效的owned入口启用。先只分配packed，完整output vector保持空。\path{program.run}等待并join全部worker，验证完成数之后，输入再无读取；最后导出只读packed及构造时缓存的n/B。此时把A的vector交换给L，逐行将上三角填为正零，再覆盖每个实际下三角位置。奇数与尾块仍截断真实维度，padded不会导出；数学和每位置k链不变。generic VM、stage、const API及结构路径继续旧分配。异常在交换前发生会释放owned A/packed，并丢弃整个M的partial。

同源control与candidate均为hybrid+快速导出+向量扫描，且均走owned API、输入销毁包含在compute；唯一编译开关差异为复用开关。418次四源FP/边界原C比较，两版分别209次，并断言每矩阵复用计数正确；另100次M4混合维度、各backend、generic读上三角与不同simd对照，18项先成功一矩阵后late-pivot/NaN/truncated失败。原A上三角包含巨大有限值、负值和负零，Cholesky最终上三角仍全部正零；generic能正确读取全部原A。API单测另外验证const输入不改、owned返回vector的allocation身份、原C整bytes。18项失败进程另记录全程peak WS/commit/CPU和宿主计量，最大工作集58,961,920字节；这些小失败Shape不代表失败N10000峰值。每项30秒退出界限，均无最终文件或partial残留。

80次资源/吞吐对照耗时201.130秒，不重算昂贵最大C oracle；所有输出整文件严格相同。下表compute三轮中位为ms，WS为该配置三轮最大工作集MiB，S/R表示同源不复用/复用。
\begin{longtable}{@{}rr rrrr@{}}
\caption{输入复用：吞吐和实际工作集独立记录}\\\toprule
$N$&worker&S compute&R compute&S WS&R WS\\\midrule\endfirsthead
\toprule$N$&worker&S compute&R compute&S WS&R WS\\\midrule\endhead
1024&1&13.528&12.444&25.6&17.6\\
1024&4&6.399&5.435&25.8&17.8\\
1024&16&10.671&9.310&26.3&18.3\\
2048&1&84.188&79.440&86.0&54.0\\
2048&4&29.370&25.657&86.2&54.2\\
2048&16&20.084&16.919&86.9&54.9\\
4096&1&672.677&599.772&326.8&198.8\\
4096&4&188.434&165.486&327.1&199.1\\
4096&16&112.035&103.354&329.6&201.6\\
10000&1&9824.830&9099.830&1924.1&1161.2\\
10000&4&3295.490&2930.670&1924.4&1161.4\\
10000&16&2152.930&2036.520&1925.8&1162.8\\
\bottomrule\end{longtable}

10000的4worker峰值工作集2,017,865,728/\allowbreak{}1,217,859,584字节，commit2,017,062,912/\allowbreak{}1,215,492,096字节，符合省掉约800,000,000字节完整L。M3=[4096,10000,4096]的16worker峰值2,019,360,768/\allowbreak{}1,219,375,104字节，不按M倍增。其4worker compute两轮中位3866.090/\allowbreak{}3639.475ms，16worker2495.990/\allowbreak{}2260.175ms。全批宿主最小可用RAM6,322,229,248字节；资源节省是直接测量，不需要将所有时间差归为缓存或分页。

增加upper清零令10000的4worker export中位33.506/\allowbreak{}57.043ms、16worker29.998/\allowbreak{}59.780ms，且这些已计入compute。依然保留单轮吞吐负结果：4worker第0轮2810.190/\allowbreak{}2930.670ms，16worker第2轮1897.410/\allowbreak{}2385.770ms；4096的16worker kernel中位74.026/\allowbreak{}78.291ms亦负。它减少整L初始分配/清零，不能据此保证每个kernel更快。这是相同hybrid布局内部的所有权单开关比较，不是相对历史row布局的单变量加速。

证据\path{artifacts/input-reuse-milestone/input-reuse-{stress,performance}-report.json}、\path{artifacts/input-reuse-lifecycle-extra/test-report.json}及\path{artifacts/input-reuse-failure-probe/test-report.json}含每个二进制/source hash、资源与失败。研究开关仍默认OFF；后续独立阶段才扩展实际结构分类失败后的auto输入复用。

\subsection{auto延迟分配与最终双配置选择}
auto扩展把结构证明放在完整L分配之前。实际带状/对角证明成功，才为结构kernel分配完整正零L，owned A仍只读；结构证明失败时，只有获完整native图证明的列块工作区可继续保持L为空，全部任务读取结束后复用A。generic、stage及const API仍保持旧分配。即使请求tiled，结构成功也优先用已验证的band/diagonal；memory配置不声称结构路径峰值也降到800MB。

同一当前源分别编译row profile和memory profile：二者owned API、扫描、B/线程策略一致；memory另外开启混合布局、转置导出及输入复用。因此这是整配置比较，不隔离所有权和布局的单一因果。418次heap/bucket四源FP压力、336次结构与新dense补充、100次backend/M4/generic原C比较，共854次完整文件比较通过。另288非有限拒绝、18多矩阵失败及24项含late-band-pivot的资源采样失败，均无输出或partial残留。

两个测试脚本问题另存证据。结构压力中曾误将M最后矩阵的\path{selected_engine}当所有矩阵的分派，错误要求混合M累计复用为零；小N可合法dense回退，runtime与原C正确。修正为aggregate计数范围、control为零、所有对角M为零，保留失败checkpoint，并不重跑已完成418项。性能脚本M拼接临时变量遮蔽报告Path，导致JSON保存至临时M oracle；长期原C参考未修改，M阶段尚未开始。主动停止自身benchmark，保存39次已逐字节通过的普通dense记录，修复变量并重建M参考，SHA与先前独立原C拼接manifest一致；再完成35次。最终74次通过，322.990秒是两段活跃耗时之和，不含中断间隙，未丢弃此前负长尾，也未重算昂贵C oracle。

下表三轮compute中位ms，WS该组三轮最高MiB；R/M为row/memory。同源数据、原C真值与每次CPU/宿主量均保存。
\begin{longtable}{@{}rr rrrr@{}}
\caption{auto整配置对照：明显吞吐负结果亦保留}\\\toprule
$N$&worker&R compute&M compute&R WS&M WS\\\midrule\endfirsthead
\toprule$N$&worker&R compute&M compute&R WS&M WS\\\midrule\endhead
2048&1&88.822&78.964&86.0&54.0\\
2048&4&33.476&27.767&86.1&54.2\\
2048&16&20.781&19.968&86.4&54.9\\
4096&1&6752.640&627.413&326.8&198.8\\
4096&4&224.084&1453.780&326.9&199.1\\
4096&16&120.411&115.141&327.2&199.7\\
10000&1&9194.470&12984.300&1924.1&1161.2\\
10000&4&3071.660&2231.550&1924.3&1161.4\\
10000&16&1729.450&1526.820&1924.5&1162.1\\
\bottomrule\end{longtable}

10000的4worker峰值WS2,017,746,944/\allowbreak{}1,217,859,584字节，commit2,016,944,128/\allowbreak{}1,215,504,384字节；M3的16worker峰值2,018,160,640/\allowbreak{}1,218,924,544字节。带状10000成功路径两者仍约1,606MB，auto compute176.412/\allowbreak{}228.659ms为负，portable204.076/\allowbreak{}202.043ms接近。M3的4worker两轮compute中位3256.665/\allowbreak{}7129.730ms明显负，范围2851.800--3661.530/\allowbreak{}3499.860--10759.600ms；16worker2221.750/\allowbreak{}1789.895ms，外部wall3366.740/\allowbreak{}3323.664ms基本接近。不能由内存节省保证吞吐改善。

10000 row的4worker第2轮18680.600ms，进程CPU68484.375ms、宿主忙55.504\%、其他CPU估计297203.125ms，可用RAM7,041,765,376字节，页错误493781次。memory单worker第1轮16008.300ms，CPU15093.750ms、其他CPU162437.500ms；M3 memory的4worker第1轮10759.600ms，CPU40046.875ms、其他CPU200750.000ms。4096 row单worker两轮6752.640/6783.120ms同样保留。全批最小可用RAM6,659,747,840字节；没有hard-fault/温度/频率计量，不把这些长尾归因于分页或某个单一DAG改变。

最终选择保守row作为\path{local-mingw}默认，另提供显式\path{memory-mingw}，输出目录\path{build-memory}。两者width16、rowgroup0、CPU检测SIMD、P/T向量化、向量扫描、选择通知开启，packed初始化正零；惰性图、等待前缀、闭式单线程、有界/锁外通知关闭。memory额外开启COLUMN/ROW/TRANSPOSE/REUSE四个选项。普通构建选项默认仍OFF，通过preset显式控制，可按资源目标选择；不以混合局部结果强行宣称全局更快。最终native优化冻结后仅执行两配置验证和报告定稿。

\subsection{最终构建与冻结验收}
两个正式preset已分别重建，native代码停止优化。\path{local-mingw}的7项CTest全部通过，159.28秒；\path{memory-mingw}亦7/7，199.42秒。每配置full基础流水线13种源、188次整文件原C比较、24项拒绝；helper49次/17拒绝，while63次，整数选择helper105次/12拒绝。因此每配置四条源码流水线共405次原C比较、53项拒绝；另有原生所有权、已知矩阵/二进制格式和308880次扫描位/保护页单测。两个构建的artifacts以build标签分开，CTest日志、结果SHA与对应binary关联保留。

最终二进制另外各做普通dense10000、真实band10000/$w=7$、M3三项，共6次整文件严格验证，46.277秒；只执行各一次，不能据其时间给稳定性能结论。默认/memory compute分别为2686.230/2454.770ms（dense）、187.096/1041.010ms（band，明显负）、3133.030/2628.060ms（M3）。dense峰值WS2,017,636,352/1,217,859,584字节，M3峰值2,017,918,976/1,218,039,808字节；band两者约1.606GB且复用计数为0、实际worker1。普通单矩阵memory复用1次，M3复用3次，default全部0。严格整文件SHA匹配既有原始C参考，结构和多矩阵边界未以小测试替代。

\path{artifacts/final-profiles/build-provenance.json}记录所有源码文件SHA、canonical source manifest SHA、preset、两版二进制和编译flags/link/cache、四条流水线与CTest日志，\path{scale-report.json}记录最后6次资源/时间/真值。preset SHA256为\path{479a0334726ca1d99d10e4f536073886e631e66174270eec564effa9ed1c571d}；源manifest为\path{f9a89e30932665402320ca7126561e07f9154a28f79ccebafd0912d32784ca94}，定义为各source相对路径到SHA字典的sort\_keys、紧凑JSON UTF8编码SHA。最终源码/编译配置快照独立保存，报告自身另有冻结SHA，不纳入此构建source manifest。

\subsection{历史整版本对照的实现限界}
历史S基线为冻结ca0版本。混合块布局为通用\path{ordered_update_avx2}增加\path{logical_size}参数，使块内stride为B时仍按实际N决定两行或四行分派。布局关闭后数学内核、归约及操作数顺序不变，但此ABI改动仍进入后续惰性图和单线程遍历构建，机器码与历史S不能假设完全相同。因此相关表格是整候选版本的同批比较，不单独证明布局、调度或计数器的因果贡献。源与构建快照保留此差异；后续扫描实验从相同当前源码、相同计时字段分别构建标量control与向量candidate，进一步减少实现混杂。所有历史逐位正确性证据仍然有效。

\subsection{同源有限值与结构扫描：局部收益和负结果}
\path{CHOL_VECTOR_SCANS}从同一当前runtime/math/ABI及计时字段构建标量control和向量candidate，减少历史构建混杂。整数位操作检查FP64指数是否全1，覆盖完整输入，含不用的上三角；每次至多读取4元素，尾部逐元素。有限值检查之后，结构扫描取每行下三角第一个非零元素，其距离等于该行最大半带宽。绝对值位掩码令$+0/-0$均为零，subnormal仍为非零。AVX2需CPU检测且\path{--simd auto}，portable仍有标量实现，数学链与融合语义不变。intrinsic位操作语义见Clang头文件文档\cite{avx2doc}。

共同压力209次、补充带宽/浮点/新稠密族336次，总计545次独立原C整文件比较通过。另288项非有限拒绝探测$\pm\mathrm{Inf}$、quiet/signaling NaN在首末、向量边界和上三角位置，均不发布输出或partial。308,880次单测与标准库标量比较，长度0--65、偏移0--7、符号零、subnormal、最大有限值及非有限位模式；Windows不可读保护页尾长0--33未越界。quick CTest两项通过。共同压力24.626秒，补充批含共同阶段76.988秒。整数位检查不对sNaN执行浮点算术；测试保证值域和拒绝行为，不承诺浮点异常标志相同。

60次交替顺序性能耗时174.734秒，复用验证过的原C真值，每次整文件核对。表中S/V为同源标量/向量，单位ms，均三轮中位；validation是compute的子集。
\begin{longtable}{@{}rr rrrr@{}}
\caption{同源扫描对照，一般稠密负结果保留}\\\toprule
$N$&worker&S compute&V compute&S validation&V validation\\\midrule\endfirsthead
\toprule$N$&worker&S compute&V compute&S validation&V validation\\\midrule\endhead
1024&1&15.998&17.065&0.308&0.238\\
1024&4&6.873&6.761&0.504&0.284\\
2048&1&109.942&151.030&1.586&1.598\\
2048&4&36.243&34.047&1.405&1.333\\
4096&1&820.546&823.971&6.591&5.172\\
4096&4&188.954&197.126&5.343&4.646\\
10000&1&9482.680&9570.510&31.473&26.005\\
10000&4&2961.660&2639.460&32.457&24.183\\
\bottomrule\end{longtable}

普通10000的4worker compute为2961.660/2639.460ms，但单worker9482.680/9570.510ms负。2048单worker109.942/151.030ms负，而validation1.586/1.598ms相近，不能归因于扫描kernel的时间差；4096的4worker188.954/197.126ms也负。输入校验几十毫秒，相对稠密$O(N^3)$数学量小，不能宣称一般稠密全局加速。

真实10000、$w=7$带状auto compute中位167.145/147.255ms，validation31.745/24.879ms，结构扫描30.885/17.077ms，外部wall1082.146/1021.630ms。请求4线程而实际带状数学仍为1worker，JSON分别记录。portable compute177.979/174.942ms近持平，结构31.114/30.313ms。完整800MB输入/输出的I/O与分配仍存在，不能把扫描阶段收益当端到端同比收益。整批最高工作集2,017,734,656字节；共享GPU/CPU负载未隔离。

\path{artifacts/vector-scans-milestone/scan-stress-report.json}和\path{scan-performance-report.json}记录补充压力、所有负值与计量，共同压力另保留\path{vector-scans-stress-report.json}。二进制与编译源/flags分别冻结。最终两个preset均显式打开扫描开关，依据有限值/结构扫描局部收益，仍保留普通dense负观测；普通CMake开关默认OFF。

\section{工作负载、硬件与测试时间成本}
\subsection{本机环境和能力需求}
本机为AMD Ryzen 9 8945HX，16物理核、32逻辑处理器；RAM为33,605,316,608字节，约31.30GiB。系统查询的L2总量为16MiB、L3总量64MiB；未设置核亲和性或隔离其他项目负载。开发环境为Windows、GCC15.2、Clang21.1、CMake4.2与Python3.12标准库venv。现有执行器使用CPU，不要求GPU或NPU。

软件能力包括C/C++编译、AST/IR与类型分析、符号循环域及依赖证明、线程同步、缓存布局、IEEE浮点语义、二进制流式IO和独立参考测试。硬件能力要求可靠FP64、足够RAM与临时磁盘；AVX2/FMA为检测后使用的加速路径，基础CPU路径保留。常规研究用4工作线程，资源伸缩研究另测8/16线程；其他项目仍可能同时开发。

\begin{table}[htbp]\centering\small
\caption{一般稠密单矩阵工作量与数据规模；十进制MB/GB}
\begin{tabular}{rrrrr}\toprule
$N$&约$N^3/3$浮点操作&单文件MB&$A+L$ MB&打包缓存约MB\\\midrule
512&44.7M&2.097&4.194&1.049\\
1024&357.9M&8.389&16.777&4.194\\
2048&2.863G&33.554&67.109&16.777\\
4096&22.906G&134.218&268.435&67.109\\
10000&333.333G&800.000&1600.000&400.000\\\bottomrule
\end{tabular}
\end{table}
公式忽略$O(N^2)$除法、开方及尾块；不是硬件执行指令计数。$N=10000,B=32$的真实打包分配为402,563,072字节，输入、输出和打包合计约2.003GB，再加任务元数据、队列、线程栈与运行环境。单矩阵的输入、参考输出、待比较输出约占2.4GB临时磁盘；流式处理使$M$影响累计时间和文件大小，而不是把全部矩阵一起驻留内存。研究参考当前保留以复用原C真值，临时混合流在核验后删除。普通/带状两组最大参考及其他里程碑已占约4.7GB，完整研究含混合M3临时输入/参考/输出宜预留至少8GB；最终清理后可只留hash与时间证据。

\subsection{计时口径}
最新图批次使用GetProcessTimes取得进程kernel与user CPU时间，合计所有线程，因此可大于外部wall；GetSystemTimes在本机32逻辑CPU上累加宿主时间，kernel含idle，计算busy时显式减idle\cite{wincpu,winprocesscpu}。宿主busy减进程CPU仅作为其他任务的近似背景，包含IO、启动及采样程序，不能分离kernel内竞争。FILETIME的单位是100ns，本机短测实际出现15.625ms量化，不把单位当精度。PageFaultCount包含所有页错误，不是hard fault或磁盘分页计数；没有性能计数器证据时不声称缓存/TLB或分页成因。
\texttt{compute\_ms}在Runtime内部围绕\texttt{execute}计时，当前CLI改用\path{execute_owned}并把输入释放计入compute；历史const入口输入通常直到写出后释放，二者不直接横比。包含输入有限值扫描、输出与打包分配、图构造、调度、线程创建/退出、kernel及工作区导出；不含文件读写。\texttt{graph\_ms}包含任务与打包缓存初始化，\texttt{kernel\_ms}包含工作线程运行与依赖队列，\texttt{export\_ms}单独记录工作区到完整L的导出，三者仍不覆盖所有compute开销。\path{input_validation_ms}单列完整输入有限值扫描，计入compute；结构扫描亦另列。\texttt{input\_read\_ms}/\texttt{output\_write\_ms}包含对应文件操作与末尾检查/刷新，Metadata读取单独计时。外部\texttt{wall\_ms}还包含进程启动、动态库加载等；不能直接把wall减compute称为纯IO。

每次性能运行之后的全文件字节比较及SHA计算不包含在Runtime compute或wall中，完整里程碑耗时包含参考编译、数据生成、C oracle、全部执行与比较。第一批采用5轮轮换顺序，第二批较大规模采用3轮或最大规模1轮；其他项目可能并发，未固定频率、冷热缓存或温度。重复读取通常使用暖文件缓存。由此可以观察局部优化收益，但不能把它们作为官方机器上的稳定得分预测。

\subsection{结构化最大Shape与一般稠密最大Shape}
早期$N=10000$测试使用完整800MB二进制布局的\emph{对角SPD}：运行时先扫描数据确认严格对角结构，再开方并保留符号零。compute为394.937ms，外部wall约6.735秒。它证明格式、最大Shape与数据驱动专门化，工作复杂度近$O(N^2)$扫描；一般稠密Cholesky约$O(N^3)$，不可用其时间替代。

本轮一般稠密$N=10000$已实际执行约333GF工作、原始C oracle及完整结果比对，约13.7秒的4线程计算也来自稠密数据。第二批整体耗时1023.768秒，主要成本包括504.670秒原始C最大规模参考、34.767秒数据生成和多次较大规模运行。后续性能搜索可复用不变源码和数据的已核验参考文件，最后仍进行完整字节核验，避免重复支付参考计算成本。

\section{本地语义限界与后续工作余量}
当前能力覆盖明确C子集与特定可证明Cholesky运算图，不能声称任意C自动并行。指针别名、二维VLA、堆分配、一般调用图、可变helper与控制流helper尚未支持。通用VM串行、受指令限额约束；有符号int32的未定义算术拒绝。Metadata是本项目受信编译器与Runtime之间的契约，Runtime不会重新从审计表达式DAG完成通用依赖证明，官方Metadata接口若另有规定仍需适配。

下一阶段价值较高的工作包括：将已验证的$O(t^2)$当前panel图扩展为一般符号依赖实例化；分散就绪队列和工作窃取以降低中心锁与优先堆成本；限制lookahead并回收已完成panel缓存；根据真实规模、线程数、ISA和成本模型选择物理$B$；证明更多合法循环和存储形式。上述改动均必须维持每位置正向归约链和独立原始C oracle。GPU、分布式或BLAS替换还需要额外浮点契约与测试，未作为已完成能力。

前端覆盖与性能优化都有剩余空间，但本地初步目标已经由实际source到动态Shape、任务DAG、严格参考结果的端到端证据支撑。正式提交前仍需对齐官方helper、评测环境、Metadata评分口径和隐藏数据；没有官方材料的部分维持显式本地假设。

\begin{table}[htbp]\centering\small
\setlength{\tabcolsep}{3pt}
\caption{后续工作余量：单人估算，非实测工期或交付承诺}
\begin{tabular}{@{}p{29mm}p{16mm}p{63mm}p{40mm}@{}}\toprule
方向&工作日&对应验证时间估算&设备与临时资源\\\midrule
前端helper/IR&5--10&更多控制流、别名、存储形式：每批5--15分钟；含新最大规模C oracle另加约8--10分钟&现有CPU；约4GB临时文件\\
稀疏专门化&4--8&不规则稀疏依赖与浮点证明：每批10--30分钟；保留带内零、FMA与完整稠密负例&现有CPU/RAM；6--10GB磁盘\\
Shape成本策略&2--4&规模、1/2/4/8线程和ISA交叉轮换：每批15--40分钟；温度/亲和性控制另计&现有CPU；无需GPU，约4GB缓存参考\\
官方集成&2--5&获得契约后格式/Metadata/机器适配：完整批次20--60分钟；隐藏评测等待不计人工&官方机器能力尚未知；本机参考约4--6GB\\\bottomrule
\end{tabular}
\end{table}
上述时间为基于已测本机成本的估算，不代表官方环境。已有实测为最新49对照加17拒绝helper批次27.294秒、计数型while批次20.073秒、结构化批次258.358秒，最大带状原始C参考另用481.611秒，一般稠密最大参考504.670秒。引用核验不变参考文件可节省重复C oracle成本；源码、输入或hash改变时须重新生成。

\appendix
\section{证据与复现}
核心文件为\texttt{tools/c\_frontend.py}、\texttt{semantic\_cholesky.py}、\texttt{analyze\_source.py}、\texttt{src/runtime.cpp}、\texttt{tiled\_runtime.cpp}与\texttt{simd\_kernel.cpp}。证据位于各里程碑JSON，报告中的数值已内嵌，阅读与编译无需访问外部数据文件。
\begin{verbatim}
.\scripts\bootstrap.ps1
.\scripts\test.ps1 -Profile full -LargeShape 10000
.venv\Scripts\python.exe tools\benchmark_tiles.py
.venv\Scripts\python.exe tools\benchmark_extended.py
.venv\Scripts\python.exe tools\benchmark_microkernels.py
.venv\Scripts\python.exe tools\benchmark_simd.py
\end{verbatim}
历史A/B工具需要保留对应基线二进制；其SHA是审计指纹，不用于计算分派。最大规模参考文件在最终核验后清理，复现长批次将重新生成数据并编译原始C参考。

本报告各性能批次使用同一官网转录C源码，SHA256为\path{0363f60d7f848b33c4e42a2be0d15d8d8d89c87caeae6b103f6a3d2059fd3a39}；左看、自设分块、显式FMA另用于正确性压力而不混入普通dense性能表。各保存的源/metadata审计及大参考manifest可交叉核查来源；较新的性能JSON直接内嵌source hash。

\begin{longtable}{@{}p{30mm}p{125mm}@{}}
\caption{历史和最终二进制SHA256}\\\toprule
版本&SHA256\\\midrule\endfirsthead
\toprule 版本&SHA256\\\midrule\endhead
阶段1旧列线程池&\scriptsize\texttt{fc69cc2f829421e4ebc022fff317a3aac6bdca0cfa00307e4b406af501701cc4}\\
紧凑图4路&\scriptsize\texttt{964b8eccdf4d95f39dfec0609587023f4c1d74c88d49e4d003a08cdeed059931}\\
候选8路&\scriptsize\texttt{8520e955958f23731a850088e3969cf776570b35364999956da64973853b7607}\\
候选16路&\scriptsize\texttt{74625dbce703bed149e4e205eeb9b30cd5a7ba61ec957ba53d45574155ed4446}\\
初代块DAG&\scriptsize\path{213fbecd3faf5ce897e71cb76a547f0b8ecc59d665a319c362933ec29888053e}\\
检测AVX2&\scriptsize\path{27057f3a7450cf52abf117f1ed0074294a12854324e1c78ec43b2b976c8a2609}\\
P/T因子缓存&\scriptsize\path{e7594cdca7a63437e753a3b7a3d9aa5e267957d04caaeb35d36c53b41d699e6f}\\
64字节对齐与heap&\scriptsize\path{665f6dc687a70cad07486c881180704897c85667ca7b9a99c3f9c9ae2b76410a}\\
两行AVX2&\scriptsize\path{cd4b421105fef36948efbb968b1b4f0c5899aec235dce2177ec4b9514266a72e}\\
两行加选择通知&\scriptsize\path{e1e6ec1bdb2de90e1e5396fc541dd50509e4c389b06150d23f13799bc249031f}\\
自动物理块候选&\scriptsize\path{03aee171a3cad76d0cfba0e6879e3406567142795cddb420981c8f94824cb6f3}\\
四行AVX2&\scriptsize\path{a69afe8a378c78c7f3e8f984edaf7b5360cac30c5d923d7f484ba37f272db954}\\
两/四行自适应&\scriptsize\path{59ff609e2c36ede9204f24ff4104680369fe5d349b3ca2ef267dac3b47123327}\\
带状路径候选&\scriptsize\path{fc511dda36fac260ef043aa4cf5a609d18c54e89a6363717b37eb97390ef2025}\\
线程相关B候选&\scriptsize\path{378179fe3efbffbb6f6d581e412a2e93613e6040acda362bcbfbd202c12ad997}\\
实验八行AVX2&\scriptsize\path{2757dd1ac400ce39aaeaa55159080129d05140728be9b8e7a710bb2013d772b7}\\
P/T SIMD未限B&\scriptsize\path{06f92cb061165868a1841182629afc134faacbcd210ff73467e4319c574bde88}\\
完整线程策略&\scriptsize\path{ca0b49fce8863690bb6cb72ffd31de1110ef2e0bb28d4e0722c95a66693fce1a}\\
有界解锁通知&\scriptsize\path{89baaba88fc8aaf2aed4a7444e1d1c0a536ad4a591ab3a9392efa74003a983ec}\\
原通知移到锁外&\scriptsize\path{25a036f5606edad86d67977f51952612ff6f4dd6440104d9c2a4a5222f3f5b10}\\
NaN缓存初始化&\scriptsize\path{2c790e79a5ada84f76ff0c7a265282f9e3daf06c20afe5677ed181c6d062099f}\\
取消缓存清零&\scriptsize\path{2c2e637e24e14faf8fabe125e517ee065e5d62ee9f1a6334797671b120478d0b}\\
连续列块工作区&\scriptsize\path{ca2feb08907f2c346cccff1e6e7193f3200b6c39389cdbad0c48e8fefcad956a}\\
混合块布局串行导出&\scriptsize\path{d721b9c855120b27c63bd7d2ee5a0a79e0bc7d75cdba343542775ed28e599fe9}\\
混合块布局快速导出&\scriptsize\path{496adfce08ac1bee8bcb79522ef9d93e74fc6cdfb2c66e7c0976993068c3b25e}\\
当前panel惰性图&\scriptsize\path{44ffaf2780f649edccf33c02ea82601ee04a16b351fa53f446d38f76afc60936}\\
当前等待与行前缀&\scriptsize\path{566326c9e80de621ec521485f3d2bcb776f28fd3064f296f00248375f18e290b}\\
单worker闭式遍历&\scriptsize\path{e81fa039c23103aefcfd954a96764c7aec2c9dfb38a88246cb67d8ff6e7d780b}\\
向量扫描标量control&\scriptsize\path{daa2b283b887d9d47ca61565afb5074072f904fd8c051d6b134505fea7a1e837}\\
向量扫描candidate&\scriptsize\path{9adb63c919737941271df4b1fefb62414e1967ace2d0cd93254c831c4b669fc8}\\
owned同源不复用control&\scriptsize\path{79ab041bc94bf4663b515819075a74511c3265d1fc1bc2dce2e68cd09fc09d54}\\
owned dense复用candidate&\scriptsize\path{6a9d3ecbf578524f3e0950f4b41e99c911de2fae10ceca46d3e75cbc86b020a6}\\
auto同源row profile&\scriptsize\path{b793c1ee2db42efd601e50233701716214e6aed3891e9df535b68978cfb05415}\\
auto memory profile&\scriptsize\path{6ee0c07b40e60bd9b38d82e8d0c8fbb4cbbd4363a4b688f182b138639b4758e6}\\
最终local-mingw默认&\scriptsize\path{2b05a6db918ced0c9c9e3bb1bfe5159ea5ed3cfee3ccb0633a538f12de409e62}\\
最终memory-mingw显式&\scriptsize\path{7da06bbd96b711f08283c1df2379741f70d8fb2be6727feae586f6c00aaad947}\\
\bottomrule
\end{longtable}

\begin{thebibliography}{99}
\bibitem{official} Chaspark，2026 Hackathon赛题19详情，2026-10-10复核。\url{https://www.chaspark.com/#/races/competitions/1326691788365586432}。
\bibitem{polly} LLVM Polly, Architecture. \url{https://polly.llvm.org/docs/Architecture.html}。
\bibitem{buttari} A. Buttari, J. Langou, J. Kurzak, J. Dongarra. A class of parallel tiled linear algebra algorithms for multicore architectures. \emph{Parallel Computing}, 35 (2009), 38--53. DOI:10.1016/j.parco.2008.10.002. 作者机构全文：\url{https://icl.utk.edu/files/publications/2010/icl-utk-440-2010.pdf}。
\bibitem{lapack} Netlib LAPACK, DPOTRF documentation. \url{https://netlib.org/lapack/explore-html/d2/d09/group__potrf_ga84e90859b02139934b166e579dd211d4.html}。
\bibitem{parsec} University of Tennessee ICL, PaRSEC Overview. \url{https://icl.utk.edu/parsec/overview/}。
\bibitem{starpu} Inria, StarPU tutorial, 2014 PATC. \url{https://starpu.gitlabpages.inria.fr/tutorials/2014-05-PATC/}。
\bibitem{llvmfp} LLVM, Language Reference: Floating-Point Environment and Fast-Math Flags. \url{https://www.llvm.org/docs/LangRef.html}。
\bibitem{clangfp} Clang, Users Manual: Floating Point Options. \url{https://clang.llvm.org/docs/UsersManual.html}。
\bibitem{repro} P. Ahrens, J. Demmel, H. D. Nguyen. Algorithms for Efficient Reproducible Floating Point Summation. \emph{ACM Transactions on Mathematical Software}, 46(3), Article 22, 2020. DOI:10.1145/3389360. 作者公开全文：\url{https://people.eecs.berkeley.edu/~demmel/ma221_Fall23/J115_Efficient_Reproducible_Summation_TOMS_2020.pdf}。
\bibitem{gccx86} GNU Compiler Collection, x86 Built-in Functions. \url{https://gcc.gnu.org/onlinedocs/gcc/x86-Built-in-Functions.html}。
\bibitem{lapackband} Netlib LAPACK, DPBTRF documentation. \url{https://www.netlib.org/lapack/explore-html/dc/d58/group__pbtrf_gafd2c3ba375ecee6dba1fba62696993e0.html}。
\bibitem{winmemory} Microsoft Learn, PROCESS\_MEMORY\_COUNTERS\_EX. \url{https://learn.microsoft.com/en-us/windows/win32/api/psapi/ns-psapi-process_memory_counters_ex}。
\bibitem{avx2doc} Clang. AVX2 intrinsic header documentation. \url{https://clang.llvm.org/doxygen/avx2intrin_8h.html}. 在线版本不等于本机实际GCC 15.2/Clang 21.1。
\bibitem{wincpu} Microsoft Learn, GetSystemTimes. \url{https://learn.microsoft.com/en-us/windows/win32/api/processthreadsapi/nf-processthreadsapi-getsystemtimes}。
\bibitem{winprocesscpu} Microsoft Learn, GetProcessTimes. \url{https://learn.microsoft.com/en-us/windows/win32/api/processthreadsapi/nf-processthreadsapi-getprocesstimes}。
\end{thebibliography}
\end{document}
